\begin{enumerate}
    \item Dans un référentiel galiléen, la variation d'énergie cinétique d'un
          solide, entre deux instants, est égale à la somme des travaux des
          forces appliquées à ce solide~:
          $\Delta E_c=E_c(B)-E_c(A)=\sum W_{AB}(\vec{F})$. \hfill\textbf{(1pt)}
    \item La trajectoire est rectiligne verticale. \hfill\textbf{(0,5pt)}
    \item On applique T.E.C~:
    \begin{gather*}
        \Delta E_c=\sum W_{AB}(\vec{F}_{\text{ext}}) \\
        \frac{1}{2}m v_B^2-\frac{1}{2}m v_A^2=W(\vec{P})=mg(z_{A}-z_{B})
        \tag{0,25} \\
        \frac{1}{2}m v^2-0=mgh \\
        \mathbox{v=\sqrt{2gh}}
        \tag{0,25}
    \end{gather*}

    Application numérique~:
    $v=\sqrt{2\times 9,8\times 15,0}$ \quad
    $\Rightarrow\ v\approx \qty{17.1}{\v}$ \hfill\textbf{(0,5pt)}
\end{enumerate}
